\documentclass{amsart}
% For dealing with references we use the comment environment
\usepackage{verbatim}
\newenvironment{reference}{\comment}{\endcomment}
%\newenvironment{reference}{}{}
\newenvironment{slogan}{\comment}{\endcomment}
\newenvironment{history}{\comment}{\endcomment}

% For commutative diagrams we use Xy-pic
\usepackage[all]{xy}

% We use 2cell for 2-commutative diagrams.
\xyoption{2cell}
\UseAllTwocells

% We use multicol for the list of chapters between chapters
\usepackage{multicol}

% This is generally recommended for better output
\usepackage{lmodern}
\usepackage[T1]{fontenc}

% For cross-file-references

% Package for hypertext links:
\usepackage{hyperref}

% For any local file, say "hello.tex" you want to link to please

% Theorem environments.
%
\theoremstyle{plain}
\newtheorem{theorem}[subsection]{Theorem}
\newtheorem{proposition}[subsection]{Proposition}
\newtheorem{lemma}[subsection]{Lemma}

\theoremstyle{definition}
\newtheorem{definition}[subsection]{Definition}
\newtheorem{example}[subsection]{Example}
\newtheorem{exercise}[subsection]{Exercise}
\newtheorem{situation}[subsection]{Situation}

\theoremstyle{remark}
\newtheorem{remark}[subsection]{Remark}
\newtheorem{remarks}[subsection]{Remarks}

\numberwithin{equation}{subsection}

% Macros
%
\def\lim{\mathop{\mathrm{lim}}\nolimits}
\def\colim{\mathop{\mathrm{colim}}\nolimits}
\def\Spec{\mathop{\mathrm{Spec}}}
\def\Hom{\mathop{\mathrm{Hom}}\nolimits}
\def\Ext{\mathop{\mathrm{Ext}}\nolimits}
\def\SheafHom{\mathop{\mathcal{H}\!\mathit{om}}\nolimits}
\def\SheafExt{\mathop{\mathcal{E}\!\mathit{xt}}\nolimits}
\def\Sch{\mathit{Sch}}
\def\Mor{\mathop{\mathrm{Mor}}\nolimits}
\def\Ob{\mathop{\mathrm{Ob}}\nolimits}
\def\Sh{\mathop{\mathit{Sh}}\nolimits}
\def\NL{\mathop{N\!L}\nolimits}
\def\CH{\mathop{\mathrm{CH}}\nolimits}
\def\proetale{{pro\text{-}\acute{e}tale}}
\def\etale{{\acute{e}tale}}
\def\QCoh{\mathit{QCoh}}
\def\Ker{\mathop{\mathrm{Ker}}}
\def\Im{\mathop{\mathrm{Im}}}
\def\Coker{\mathop{\mathrm{Coker}}}
\def\Coim{\mathop{\mathrm{Coim}}}

% Boxtimes
%
\DeclareMathSymbol{\boxtimes}{\mathbin}{AMSa}{"02}

%
% Macros for moduli stacks/spaces
%
\def\QCohstack{\mathcal{QC}\!\mathit{oh}}
\def\Cohstack{\mathcal{C}\!\mathit{oh}}
\def\Spacesstack{\mathcal{S}\!\mathit{paces}}
\def\Quotfunctor{\mathrm{Quot}}
\def\Hilbfunctor{\mathrm{Hilb}}
\def\Curvesstack{\mathcal{C}\!\mathit{urves}}
\def\Polarizedstack{\mathcal{P}\!\mathit{olarized}}
\def\Complexesstack{\mathcal{C}\!\mathit{omplexes}}
% \Pic is the operator that assigns to X its picard group, usage \Pic(X)
% \Picardstack_{X/B} denotes the Picard stack of X over B
% \Picardfunctor_{X/B} denotes the Picard functor of X over B
\def\Pic{\mathop{\mathrm{Pic}}\nolimits}
\def\Picardstack{\mathcal{P}\!\mathit{ic}}
\def\Picardfunctor{\mathrm{Pic}}
\def\Deformationcategory{\mathcal{D}\!\mathit{ef}}

\setlength{\emergencystretch}{3em}
\begin{document}
\title[Stacks Project, Tag 0E25]{491 --- Scheme pushouts along closed immersions and integral morphisms exist under the stated affine-neighbourhood compatibility hypothesis}
\maketitle
\noindent Copyright (C) 2005--2025 Johan de Jong. Stacks Project authors. Source: \href{https://stacks.math.columbia.edu/tag/0E25}{Tag 0E25}.

\noindent Editorial difficulty note (not author text). Difficulty H2: The topological pushout with its fibre-product structure sheaf must be proved locally affine at every glued point. Compatible affine neighbourhoods reduce this to a ring fibre product whose localizations, integral projection and tensor-product identity are checked, and the local pushout universal properties must then descend globally..

\noindent Editorial provenance note (not author text). History: excerpt from revision \texttt{a0444\allowbreak 6e57e\allowbreak c1fbc\allowbreak 252a8\allowbreak 71afc\allowbreak ec775\allowbreak 2fb28\allowbreak 07b14}, more-morphisms.tex, lines 19757--19844. Reference presentation was adapted; original-author context and core support blocks were retained or added. The mathematical wording is unchanged. Licensed under GNU FDL 1.2 or later; no invariant sections or cover texts. See ../licenses/COPYING. Context blocks reproduce author definitions and supporting results; they are not additional counted problems.

\begin{situation}
\label{situation-pushout-along-closed-immersion-and-integral}
Here $S$ is a scheme and $i : Z \to X$ and $j : Z \to Y$
are morphisms of schemes over $S$. We assume
\begin{enumerate}
\item $i$ is a closed immersion,
\item $j$ is an integral morphism of schemes,
\item for $y \in Y$ there exists an affine open $U \subset X$
with $j^{-1}(\{y\}) \subset i^{-1}(U)$.
\end{enumerate}
\end{situation}

\begin{proposition}
\label{proposition-pushout-along-closed-immersion-and-integral}
\begin{reference}
\href{https://stacks.math.columbia.edu/bibliography}{Bibliography: Ferrand-Conducteur (Theorem 7.1 part iii)}
\end{reference}
In Situation \ref{situation-pushout-along-closed-immersion-and-integral}
the pushout $Y \amalg_Z X$ exists in the category of schemes. Picture
$$
\xymatrix{
Z \ar[r]_i \ar[d]_j & X \ar[d]^a \\
Y \ar[r]^-b & Y \amalg_Z X
}
$$
The diagram is a fibre square, the morphism $a$ is integral,
the morphism $b$ is a closed immersion, and
$$
\mathcal{O}_{Y \amalg_Z X} =
b_*\mathcal{O}_Y \times_{c_*\mathcal{O}_Z} a_*\mathcal{O}_X
$$
as sheaves of rings where $c = a \circ i = b \circ j$.
\end{proposition}

\begin{proof}
As a topological space we set $Y \amalg_Z X$ equal to the pushout of
the diagram in the category of topological spaces (Topology, Section
\href{https://stacks.math.columbia.edu/tag/0B1W}{section colimits (Tag 0B1W)}). This is just the pushout
of the underlying sets (Topology, Lemma \href{https://stacks.math.columbia.edu/tag/0B1X}{lemma colimits (Tag 0B1X)})
endowed with the quotient topology.
On $Y \amalg_Z X$ we have the maps of sheaves of rings
$$
b_*\mathcal{O}_Y \longrightarrow c_*\mathcal{O}_Z
\longleftarrow a_*\mathcal{O}_X
$$
and we can define
$$
\mathcal{O}_{Y \amalg_Z X} =
b_*\mathcal{O}_Y \times_{c_*\mathcal{O}_Z} a_*\mathcal{O}_X
$$
as the fibre product in the category of sheaves of rings. To prove that we
obtain a scheme we have to show that every point has an
affine open neighbourhood. This is clear for points not in the image of $c$
as the image of $c$ is a closed subset whose complement is
isomorphic as a ringed space to $(Y \setminus j(Z)) \amalg (X \setminus i(Z))$.

\medskip\noindent
A point in the image of $c$ corresponds to a unique $y \in Y$
in the image of $j$. By
Lemma \href{https://stacks.math.columbia.edu/tag/0ECJ}{lemma prepare pushout along closed immersion and integral (Tag 0ECJ)}
we find affine opens $U \subset X$ and $V \subset Y$ with
$y \in V$ and $i^{-1}(U) = j^{-1}(V)$.
Since the construction of the first paragraph is clearly compatible
with restriction to compatible open subschemes, to prove that it
produces a scheme we may assume $X$, $Y$, and $Z$ are affine.

\medskip\noindent
If $X = \Spec(A)$, $Y = \Spec(B)$, and $Z = \Spec(C)$ are affine, then
More on Algebra, Lemma \href{https://stacks.math.columbia.edu/tag/0B7J}{lemma points of fibre product (Tag 0B7J)}
shows that $Y \amalg_Z X = \Spec(B \times_C A)$ as topological spaces.
To finish the proof that $Y \times_Z X$ is a scheme, it suffices to show
that on $\Spec(B \times_C A)$ the structure sheaf is the fibre product
of the pushforwards. This follows by applying
More on Algebra, Lemma \href{https://stacks.math.columbia.edu/tag/01Z8}{lemma diagram localize (Tag 01Z8)}
to principal affine opens of $\Spec(B \times_C A)$.

\medskip\noindent
The discussion above shows the scheme $Y \amalg_Z X$
has an affine open covering $Y \amalg_Z X = \bigcup W_i$
such that $U_i = a^{-1}(W_i)$, $V_i = b^{-1}(W_i)$, and
$\Omega_i = c^{-1}(W_i)$ are affine open in $X$, $Y$, and $Z$.
Thus $a$ and $b$ are affine.
Moreover, if $A_i$, $B_i$, $C_i$ are the rings corresponding to
$U_i$, $V_i$, $\Omega_i$, then $A_i \to C_i$ is surjective and
$W_i$ corresponds to $A_i \times_{C_i} B_i$ which surjects onto
$B_i$. Hence $b$ is a closed immersion.
The ring map $A_i \times_{C_i} B_i \to A_i$ is integral by
More on Algebra, Lemma \href{https://stacks.math.columbia.edu/tag/0E1S}{lemma fibre product integral (Tag 0E1S)}
hence $a$ is integral. The diagram is cartesian because
$$
C_i \cong B_i \otimes_{B_i \times_{C_i} A_i} A_i
$$
This follows as $B_i \times_{C_i} A_i \to B_i$
and $A_i \to C_i$ are surjective maps whose kernels are the same.

\medskip\noindent
Finally, we can apply Lemmas \href{https://stacks.math.columbia.edu/tag/0ET0}{lemma basic example pushout (Tag 0ET0)} and
\href{https://stacks.math.columbia.edu/tag/0BMP}{lemma pushout fpqc local (Tag 0BMP)} to conclude our construction is a pushout
in the category of schemes.
\end{proof}
\end{document}
